% =============================================================================
% SECTION 3 — METHOD
% =============================================================================
% Drop-in replacement for the Method section of paper_skeleton.tex.
% Uses standard amsmath / algorithm / algpseudocode environments.
%
% Required packages already in skeleton:
%   amsmath, algorithm, algpseudocode, multirow, booktabs, graphicx
%
% Forward references used:
%   Section~\ref{sec:exp:ablation}  (selection-algorithm ablation)
%   Section~\ref{sec:method:cot}     (CoT augmentation, this section)
%   Appendix~\ref{app:algo}          (full algorithm details)
% =============================================================================

\section{Method}
\label{sec:method}

\subsection{Motivation: Function Calls as Agent-Like Structures}
\label{sec:method:motivation}

A coding agent operates by interleaving \emph{actions} (tool calls) with
\emph{observations} (tool returns) and continuing its reasoning conditioned on
both. Formally, at step $t$ an agent maintains a history $h_t$ and samples an
action $a_t \sim \pi(a_t \mid h_t)$; the environment yields an observation
$o_{t+1} \sim p(o_{t+1} \mid h_t, a_t)$; and the agent then produces a
continuation conditioned on the entire trace
$\pi(\,\cdot \mid h_t, a_t, o_{t+1})$. The capability that distinguishes a
useful agent from a competent code model is precisely this third step:
\emph{integrating an externally returned value into ongoing reasoning}.

We observe that this conditioning structure is not unique to agentic
trajectories. Every function call site in real-world code exhibits the same
four-part pattern:
\textit{(i)} a pre-call context that establishes intent and binds arguments,
\textit{(ii)} the call itself,
\textit{(iii)} a return value produced by code outside the immediate scope,
and \textit{(iv)} downstream code that consumes that return value. The pre-call
context corresponds to $h_t$, the call to $a_t$, the return to $o_{t+1}$, and
the downstream usage to the agent's continuation. Figure~\ref{fig:isomorphism}
illustrates this correspondence.

% TODO FIG 2: side-by-side isomorphism diagram
% \begin{figure}[t]
% \centering
% \includegraphics[width=\linewidth]{figs/isomorphism.pdf}
% \caption{Structural correspondence between a function call site (top) and a
% single agent step (bottom). Both decompose into context $\to$ action $\to$
% external return $\to$ continuation. We exploit this correspondence by
% constructing a self-supervised mid-training objective from real-world code.}
% \label{fig:isomorphism}
% \end{figure}

This correspondence motivates our central design: rather than rely on
synthesised agent trajectories, we extract an agentic conditioning signal from
the structure already present in source code, using the
fill-in-the-middle (FIM) objective. We emphasise that the correspondence is
\emph{structural} rather than literal. FIM training maximises
$\log P(M \mid \text{prefix}, \text{suffix})$, where the suffix is given;
agent inference generates the continuation autoregressively. We hypothesise
that training a model to reason bidirectionally about function-level
dependencies induces a representation that is also useful when the
``suffix'' is generated rather than observed---an empirical question we test
in Section~\ref{sec:exp}.

Standard random-span FIM, as used during pretraining of recent code
LLMs~\citep{TODO_qwen_coder, TODO_deepseek_coder, TODO_starcoder}, captures
this structure only incidentally: most masked spans cut across syntactically
arbitrary boundaries and provide weak agent-relevant signal. We instead
construct a \emph{function-aware} variant that selects masking targets
based on program structure and contextual predictability.

\subsection{Data Collection and Decontamination}
\label{sec:method:data}

We curate a corpus of $990$ Python repositories from GitHub. Starting from
$\sim\!2{,}000$ candidates retrieved by combining a star-count threshold with
ten topic queries (Appendix~\ref{app:data}), we apply manual quality filtering
and remove every repository whose origin overlaps with the source repositories
of SWE-Bench~\citep{TODO_swebench} (verified by repository name and any known
fork). To eliminate the possibility of test-time leakage from training-data
commits, we restrict each repository to commits whose timestamp precedes the
earliest base-commit used in SWE-Bench-Verified and SWE-Bench-Lite. The full
repository list, license inventory, and commit hashes will be released to
support reproduction.

% TODO FIG 3: 10-category distribution
% TODO TAB 1: dataset statistics (file count, total LoC, license breakdown)

\subsection{Function-Aware FIM Target Selection}
\label{sec:method:selection}

Given the file source of each repository, our pipeline produces a set of
mask targets together with the corresponding masked file. Each target is
a single function or a connected group of $2$--$3$ functions; the latter
variant is studied in Section~\ref{sec:exp:ablation}. The pipeline has four
stages: dependency-graph construction, complexity scoring, inferability
scoring, and threshold-based selection. Algorithm~\ref{alg:selection}
summarises the single-function pipeline; group selection is described
in Section~\ref{sec:method:groups}.

\subsubsection{Program Dependency Graph}
\label{sec:method:pdg}

For each source file we parse its abstract syntax tree (AST) and extract
the set $\mathcal{V}$ of function nodes, where each node corresponds to a
top-level function or a class method identified by its qualified name
(e.g.\ \texttt{ClassName.method}). We construct two edge sets:

\paragraph{Call edges $\mathcal{E}_{\text{call}}$.}
For every \texttt{Call} node inside the body of $u \in \mathcal{V}$, we
attempt to resolve its callee to some $v \in \mathcal{V}$, adding a directed
edge $u \to v$. Resolution handles three Python idioms:
direct module-level invocations, class instantiation
(\texttt{ClassName(...)} resolved to \texttt{ClassName.\_\_init\_\_}), and
method calls of the form \texttt{self.m(...)} or \texttt{cls.m(...)}
inside class bodies. When a call cannot be resolved exactly we fall back to
short-name matching against the qualified-name index, which recovers most
cross-module references while admitting a controlled false-positive rate
(see Appendix~\ref{app:algo}).

\paragraph{Sibling edges $\mathcal{E}_{\text{sib}}$.}
For every pair of methods belonging to the same class we add an undirected
sibling edge. These edges capture intra-class coupling that does not
manifest as direct calls but does carry information through shared
instance state.

\subsubsection{Complexity Score $\hat{H}$}
\label{sec:method:complexity}

A useful FIM target should require non-trivial reasoning to reconstruct.
For each $v \in \mathcal{V}$ we define
%
\begin{equation}
  \hat{H}(v) \;=\; w_{\ell}\,\phi\!\left(\mathrm{LoC}(v),\,c_{\ell}\right)
            + w_{c}\,\phi\!\left(\mathrm{CC}(v),\,c_{c}\right)
            + w_{d}\,\phi\!\left(\mathrm{D}(v),\,c_{d}\right),
  \label{eq:complexity}
\end{equation}
%
where $\mathrm{LoC}(v)$ is the lines-of-code count, $\mathrm{CC}(v)$ is
McCabe cyclomatic complexity, $\mathrm{D}(v)$ is the maximum nesting depth
of control-flow constructs, and $\phi(x, c) = \min(x / c,\, 2)$ normalises
each quantity by a soft cap. We use $(c_{\ell}, c_{c}, c_{d}) = (50, 10, 5)$
and $(w_{\ell}, w_{c}, w_{d}) = (0.4, 0.4, 0.2)$. The cap at $2$ rather than
$1$ allows genuinely complex functions to score above the median without
letting outliers dominate.

\subsubsection{Inferability Score $\hat{I}$}
\label{sec:method:inferability}

A target should also be \emph{recoverable} from the surrounding context.
$\hat{I}(v)$ aggregates five context-derived signals that approximate the
mutual information between $v$'s body and the rest of the file:
%
\begin{equation}
  \hat{I}(v) \;=\; \alpha\,C_{\mathrm{caller}}(v)
              + \beta\,C_{\mathrm{callee}}(v)
              + \gamma\,C_{\mathrm{sig}}(v)
              + \delta\,C_{\mathrm{doc}}(v)
              + \varepsilon\,C_{\mathrm{class}}(v),
  \label{eq:inferability}
\end{equation}
%
with $(\alpha,\beta,\gamma,\delta,\varepsilon) = (0.30, 0.25, 0.20, 0.10, 0.15)$.

\paragraph{Caller signal $C_{\mathrm{caller}}$.}
For each caller $u$ of $v$, we score the specificity of its call sites by
inspecting the syntactic form of each argument: literal constants are
maximally informative, named variables less so, keyword arguments contribute
proportionally, and a base score is awarded for the call site itself. The
sum is normalised by an expected-callers constant.

\paragraph{Callee signal $C_{\mathrm{callee}}$.}
The number of intra-file functions that $v$ itself calls, normalised by a
small expected-fan-out constant. A function that orchestrates known
helpers is more recoverable than one that calls only opaque externals.

\paragraph{Signature signal $C_{\mathrm{sig}}$.}
A monotone aggregate of return-type annotation, parameter-type annotations,
the descriptiveness of the function name (decomposed into snake-case parts),
and the count of non-\texttt{self} parameters.

\paragraph{Documentation signal $C_{\mathrm{doc}}$.}
Indicator of docstring presence. We do not attempt to score docstring
informativeness, both for robustness and because Section~\ref{sec:method:cot}
generates a richer textual rationale.

\paragraph{Class signal $C_{\mathrm{class}}$.}
For methods, the number of in-class siblings (normalised) plus a small bonus
when an \texttt{\_\_init\_\_} is present in the same class.

We deliberately keep each component a hand-designed proxy rather than a
learned predictability score. A learned proxy would couple selection to a
particular reference model and complicate the cross-base-model
generalisation analysis (Section~\ref{sec:exp:crossbase}); hand-designed
features yield a model-agnostic curriculum.

\subsubsection{Single-Function Score}
\label{sec:method:single}

The final score combines $\hat{H}$ and $\hat{I}$ in a form that rewards
\emph{simultaneously} high values of both:
%
\begin{equation}
  \mathrm{FIM}(v) \;=\;
  \underbrace{\frac{\hat{H}(v)\,\hat{I}(v)}{\hat{H}(v) + \hat{I}(v) + \epsilon}}_{\text{base score}}
  \;\cdot\;
  \underbrace{\rho\!\left(\Delta(v)\right)}_{\text{difficulty penalty}},
  \label{eq:fim_single}
\end{equation}
%
where $\epsilon$ is a small constant for numerical stability. The base score
is a harmonic-mean-like combination: a pure product
$\hat{H}\,\hat{I}$ would treat $(\hat{H}, \hat{I}) = (\text{high},
\text{high})$ as comparable to $(\text{very high}, \text{low})$, whereas
Eq.~\eqref{eq:fim_single} forces both quantities to be large simultaneously
and penalises imbalance.

\paragraph{Difficulty and the one-sided Gaussian penalty.}
We define the residual difficulty as the share of complexity that is not
explained by context:
\begin{equation}
  \Delta(v) \;=\; \frac{\max\!\left(0,\,\hat{H}(v) - \hat{I}(v)\right)}{\hat{H}(v) + \epsilon}
  \;\in\; [0, 1).
\end{equation}
The penalty $\rho$ is one-sided: targets with $\Delta(v) \le \tau_d$ are
left intact, while targets above the threshold are damped:
%
\begin{equation}
  \rho(\Delta) \;=\;
  \begin{cases}
    1, & \Delta \le \tau_d, \\[2pt]
    \exp\!\left(-\dfrac{(\Delta - \tau_d)^2}{2\sigma^2}\right), & \Delta > \tau_d.
  \end{cases}
\end{equation}
We use $\tau_d = 0.5$ and $\sigma = 0.2$. The asymmetry reflects an asymmetry
in the underlying objective: trivially easy targets are eliminated by the
hard floor on $\hat{H}$ (below), so we do not need to additionally reward
high inferability; conversely, targets that remain hard even with full
context become unlearnable noise and must be down-weighted.

\paragraph{Hard filters.}
Independent of the smooth score, we discard
\textit{(i)} functions with $\mathrm{LoC}(v) \notin [10, 200]$ to control
training-instance length;
\textit{(ii)} dunder methods (\texttt{\_\_init\_\_}, \texttt{\_\_repr\_\_},
etc.), which serve as scaffolding rather than reasoning targets;
\textit{(iii)} functions with $\hat{H}(v) < 0.15$ to remove trivial cases;
and \textit{(iv)} functions with $\mathrm{FIM}(v) < 0.08$.
A file with no surviving target contributes nothing to the corpus, ensuring
each training sample carries informative signal.

\begin{algorithm}[t]
\caption{Single-function FIM target selection for one file.}
\label{alg:selection}
\begin{algorithmic}[1]
  \Require Source file $s$; thresholds $\tau_d, \tau_{\mathrm{FIM}}, \tau_{\hat H}$
  \State $T \gets \mathrm{ParseAST}(s)$
  \State $\mathcal{V}, \mathcal{E}_{\mathrm{call}}, \mathcal{E}_{\mathrm{sib}} \gets \mathrm{BuildPDG}(T)$
        \Comment{Section~\ref{sec:method:pdg}}
  \State $\mathcal{C} \gets \emptyset$
  \For{each $v \in \mathcal{V}$}
    \If{$v$ is a dunder method \textbf{or} $\mathrm{LoC}(v) \notin [10, 200]$}
      \State \textbf{continue}
    \EndIf
    \State Compute $\hat{H}(v)$ and $\hat{I}(v)$
          \Comment{Eqs.~\eqref{eq:complexity}, \eqref{eq:inferability}}
    \If{$\hat{H}(v) < \tau_{\hat H}$} \textbf{continue} \EndIf
    \State $\Delta(v) \gets \max(0, \hat{H}(v)-\hat{I}(v)) / (\hat{H}(v)+\epsilon)$
    \State $\mathrm{FIM}(v) \gets \frac{\hat{H}(v)\,\hat{I}(v)}{\hat{H}(v)+\hat{I}(v)+\epsilon}\cdot \rho(\Delta(v))$
    \If{$\mathrm{FIM}(v) \ge \tau_{\mathrm{FIM}}$}
      \State $\mathcal{C} \gets \mathcal{C} \cup \{v\}$
    \EndIf
  \EndFor
  \State \Return $\mathcal{C}$ \Comment{ranked by $\mathrm{FIM}$ descending}
\end{algorithmic}
\end{algorithm}

\subsubsection{Multi-Function Group Selection}
\label{sec:method:groups}

Real-world code patches frequently span multiple related
functions~\citep{swebench}, motivating an extension that masks
\emph{groups} of $k = 2$ or $k = 3$ structurally connected functions.
We use this variant in the granularity ablation
(Section~\ref{sec:exp:ablation}); single-function masking remains the
default in our main results.

\paragraph{Group topology.}
A candidate group must form a connected subgraph in
$\mathcal{E}_{\mathrm{call}} \cup \mathcal{E}_{\mathrm{sib}}$. We enumerate
eight topology patterns:
\textit{caller-callee}, \textit{co-callee} (two callees of a common caller),
\textit{sibling-coupled} (same-class methods sharing $\geq 1$ instance
attribute), and \textit{mutual-call} for $k=2$;
\textit{call-chain} ($A \!\to\! B \!\to\! C$),
\textit{hub} ($A \!\to\! B,\; A \!\to\! C$),
\textit{fan-in} ($B \!\to\! A,\; C \!\to\! A$),
and \textit{class-triad} (three pairwise state-sharing methods of the same
class) for $k=3$. The full taxonomy is detailed in Appendix~\ref{app:algo}.

\paragraph{Group score.}
For a group $G = \{v_1, \dots, v_k\}$ we define
%
\begin{equation}
  \mathrm{FIM}(G) \;=\;
  \underbrace{\mathrm{Coup}(G)}_{\text{coupling}}
  \;\cdot\;
  \frac{\hat{H}(G)\,\hat{I}(G)}{\hat{H}(G) + \hat{I}(G) + \epsilon}
  \;\cdot\;
  \rho\!\left(\Delta(G)\right),
  \label{eq:fim_group}
\end{equation}
%
where $\hat{H}(G) = \tfrac{1}{k}\sum_v \hat{H}(v)$ averages individual
complexities, $\Delta(G)$ is defined analogously to the single-function case,
and $\mathrm{Coup}(G) \in [0, 1]$ is a normalised count of intra-group
edges plus a shared-state ratio:
%
\begin{equation}
  \mathrm{Coup}(G)
  \;=\; w_{\mathrm{c}}\,\frac{|E^{\text{call}}_G|}{k(k\!-\!1)}
   + w_{\mathrm{s}}\,\frac{|E^{\text{sib}}_G|}{\binom{k}{2}}
   + w_{\mathrm{st}}\,\overline{\mathrm{Jacc}}_G,
\end{equation}
where $\overline{\mathrm{Jacc}}_G$ is the mean Jaccard similarity over
pairs in $G$ between sets of accessed instance attributes
(\texttt{self.x}). We use $(w_c, w_s, w_{\mathrm{st}}) = (0.50, 0.20, 0.30)$.
The coupling term enters multiplicatively because, unlike the
single-function case, the value of a multi-function target derives
specifically from cross-function dependencies; a weakly coupled group is
indistinguishable from $k$ independent samples.

\paragraph{Group inferability.}
$\hat{I}(G)$ is recomputed under the assumption that \emph{all} members of
$G$ are masked simultaneously: caller, callee, and sibling contributions
from inside $G$ are excluded, and docstrings of group members
contribute zero (they are part of the masked body). Signature information
is retained because we do not mask function signatures. This re-computation
prevents groups from receiving spurious credit for intra-group references
that disappear under joint masking.

\paragraph{Group selection.}
We score every candidate group emitted by the topology enumerator and
discard those with insufficient coupling, complexity, or score (defaults
in Appendix~\ref{app:algo}). Groups exceeding LoC ratios of
$0.30$ ($k=2$) or $0.40$ ($k=3$) of the file are rejected. The
remaining candidates are sorted by $\mathrm{FIM}(G)$ and selected greedily
under non-overlap: a function may belong to at most one selected group per
file. Per-file caps ($5$ pairs, $3$ triples) prevent any single file from
dominating the corpus.

\subsection{Chain-of-Thought Augmentation}
\label{sec:method:cot}

For each selected target we generate a natural-language rationale that
explicates how the function's behaviour is recoverable from the surrounding
file. The rationale is produced by Gemini-3-Preview prompted with the
masked file and the original function body (the prompt template is in
Appendix~\ref{app:algo}). At training time, the rationale is embedded
\emph{inside} the FIM middle span, immediately before the function body:
%
\begin{center}\small
\texttt{<fim\_prefix>}\,$\langle$prefix$\rangle$\,%
\texttt{<fim\_suffix>}\,$\langle$suffix$\rangle$\,%
\texttt{<fim\_middle>}\,$\langle$rationale$\rangle$\,$\langle$body$\rangle$
\end{center}
%
The model is therefore trained to produce, in order, a reasoning trace
followed by the code consistent with that trace---mirroring the
think-then-act structure of an agent step.

We emphasise that the role of the rationale is not to distil the teacher's
generic capability, but to provide a structured intermediate supervision
that aligns reasoning with code under the FIM objective. To separate the
contribution of FIM structure from that of CoT distillation, we ablate
both the presence and the source of the rationale in
Section~\ref{sec:exp:ablation}, including a \emph{self-CoT} variant in
which the same Qwen2.5-Coder-Instruct model that we are training
generates its own rationales. We report that the FIM structure alone
(no CoT) already accounts for a substantial fraction of the gain, and that
self-generated CoT recovers most of the remainder---indicating that
high-quality rationales help but are not the principal driver.
% NOTE: 这一段的具体数字在 ablation 数据出来后填入

\subsection{Mid-Training and Post-Training Setup}
\label{sec:method:setup}

We mid-train Qwen2.5-Coder-Instruct (7B, 14B, 32B) and Qwen3-8B-Base on the
selected FIM corpus. The training objective is the standard FIM loss on
the middle span (rationale plus body), packed to the model's native
context length and using the model's native FIM sentinel tokens to
preserve compatibility with any FIM signal already learned during
pretraining. After mid-training we apply existing agentic
post-training pipelines (r2e-gym~\citep{TODO_r2egym} and
swe-smith~\citep{TODO_swesmith}) without modification. Hyperparameters,
token budgets, and learning-rate schedules are reported in
Appendix~\ref{app:hp}.